\documentclass[10pt,letterpaper]{article}
\usepackage[lmargin=2.2in,rmargin=0.7in,tmargin=1.0in,bmargin=1.0in]{geometry}
\usepackage{notes}
\usetikzlibrary{decorations.pathmorphing}

\begin{document}

\begin{center}
\Large\textbf{2.4 -- Characteristic Equation: Complex Roots}\\
\normalsize (Section 2.2.3)
\end{center}

%%%%%%%%%%%%%%%%%%%%%%%%%%%%%%%%%%%%%%%%%%%%%%%%%%%%%%%%%%%%%%%%%%%%%%%%
\section{Theory of Complex Solutions}

\subsection{Complex Eigenvalues}

Last lecture, we found the eigenvalues of
\[
x''-2x'+5x=0.
\]
The characteristic equation gave us two complex eigenvalues
\[
\lambda_1=1+2i,\quad \lambda_2=1-2i.
\]
We can write them both compactly using the $\pm$ symbol
$$
\lambda = 1\pm 2i.
$$
The complex eigenvalues always come in pairs in the form
\[
\lambda=\alpha\pm i\beta,
\qquad \alpha\in\mathbb{R},\quad \beta>0.
\]
Just as before, each eigenvalue gives an exponential solution.
The next step is to use these complex solutions to obtain
real solutions.

%-------------------------------------------------------------
\subsection{Using Euler's Formula}

The eigenvalue $\lambda=1+2i$ gives the complex solution
\[
z(t)=e^{(1+2i)t}.
\]
Recall Euler's formula from the last lecture:
\[
e^{it}=\cos t+i\sin t.
\]
Use it to separate the real and imaginary parts:
\begin{align*}
z(t)
&=e^t e^{2it}\\
&=e^t(\cos2t+i\sin2t)\\
&=\underbrace{e^t\cos2t}_{\text{real}}
+i\underbrace{e^t\sin2t}_{\text{imaginary}}.
\end{align*}
These two real functions turn out to be solutions to our DE.
The next theorem explains why.
%-------------------------------------------------------------
\subsection{Theorem about Complex Solutions}

\begin{theorem}
Suppose $z(t)=f(t)+ig(t)$ is a complex solution to
\[
az''+bz'+cz=0,
\]
where $a,b,c$ are real constants and $a\ne0$, and $f,g$ are
real-valued functions. Then $f(t)$ and $g(t)$ are also
solutions to the same DE.
\end{theorem}

\begin{proof}
Substitute $z=f+ig$ into the DE:
\begin{align*}
0
&=az''+bz'+cz\\
&=a(f''+ig'')+b(f'+ig')+c(f+ig)\\
&=(af''+bf'+cf)+i(ag''+bg'+cg).
\end{align*}
Two complex numbers are equal when their real parts match
and their imaginary parts match. Comparing with $0=0+0i$ gives
\[
af''+bf'+cf=0,
\qquad
ag''+bg'+cg=0.
\]
Thus, $f$ and $g$ each satisfy the original DE.
\end{proof}

This theorem gives us two real solutions from a single complex
solution. In practice, we choose one complex eigenvalue, write
$e^{\lambda t}$ using Euler's formula, and take its real and
imaginary parts. For the complex roots considered here, these
two functions are linearly independent, so their linear
combination gives the real general solution. We are now ready for an example.

%-------------------------------------------------------------
\begin{example}[1]
{
Find the real general solution to
\[
x''-2x'+5x=0.
\]
}
We already found the eigenvalues $\lambda=1\pm2i$.
Using $\lambda=1+2i$, we obtain the complex solution
\[
z(t)=e^{(1+2i)t}.
\]
Apply Euler's formula to identify its real and imaginary parts:
\begin{align*}
z(t)
&=e^t e^{2it}\\
&=e^t(\cos2t+i\sin2t)\\
&=\underbrace{e^t\cos2t}_{\text{real}}
+i\underbrace{e^t\sin2t}_{\text{imaginary}}.
\end{align*}
By the theorem, the functions
\[
x_1(t)=e^t\cos2t,
\qquad
x_2(t)=e^t\sin2t
\]
are real solutions. They are linearly independent, so the real
general solution is
\[
\boxed{
x(t)=C_1e^t\cos2t+C_2e^t\sin2t.
}
\]
\end{example}

%%%%%%%%%%%%%%%%%%%%%%%%%%%%%%%%%%%%%%%%%%%%%%%%%%%%%%%%%%%%%%%%%%%%%%%%
\section{IVPs}

\subsection{General Solution Formula}

Suppose the characteristic equation
\[
a\lambda^2+b\lambda+c=0
\]
has a negative discriminant $D=b^2-4ac<0$. Its roots have the form
\[
\lambda=\alpha\pm i\beta,
\qquad \beta>0.
\]
Euler's formula gives
\[
e^{(\alpha+i\beta)t}
=e^{\alpha t}\cos(\beta t)
+i e^{\alpha t}\sin(\beta t).
\]
The real and imaginary parts form two linearly independent
solutions. Therefore, the real general solution is
\[
\boxed{
x(t)=e^{\alpha t}
\left(C_1\cos(\beta t)+C_2\sin(\beta t)\right).
}
\]
You should commit this formula to memory. Identify $\alpha$
and $\beta$ from the eigenvalues, write the general solution,
and then use the initial conditions to find $C_1$ and $C_2$. If you happen to forget the general solution formula, you can use 
Euler's formula to derive the equation.

%-------------------------------------------------------------
\begin{example}[2]
{
Find the solution to the IVP
\[
\begin{dcases}
x''+2x'+10x=0,\\
x(0)=1,\\
x'(0)=5.
\end{dcases}
\]
}
The characteristic equation is
\[
\lambda^2+2\lambda+10=0.
\]
We could use the quadratic formula, but let's \textbf{complete the square} instead:
\begin{align*}
\lambda^2+2\lambda&=-10 &&\com{move 10 to RHS}\\
\lambda^2+2\lambda+1&=-9 &&\com{add $(2/2)^2$ to both sides}\\
(\lambda+1)^2&=-9  &&\com{rewrite LHS as the square}\\
\lambda+1&=\pm3i &&\com{square root both sides}\\
\lambda&=-1\pm3i &&\com{subtract 1}.
\end{align*}
Thus, $\alpha=-1$ and $\beta=3$. The general solution is
\[
x(t)=e^{-t}(C_1\cos3t+C_2\sin3t).
\]
Apply the first initial condition:
\[
x(0)=e^0(C_1\cos0+C_2\sin0)=C_1\equiv1.
\]
With $C_1=1$, differentiate using the product rule:
\begin{align*}
x'(t)
&=-e^{-t}(\cos3t+C_2\sin3t)\\
&\qquad+e^{-t}(-3\sin3t+3C_2\cos3t).
\end{align*}
The second initial condition gives
\[
x'(0)=-1+3C_2\equiv5
\quad\Rightarrow\quad C_2=2.
\]
Therefore,
\[
\boxed{
x(t)=e^{-t}(\cos3t+2\sin3t).
}
\]
\end{example}

% ------------------

\subsection{Solution to IVP}
Suppose you are given the solution of the form 
$$
x(t)=e^{\alpha t}
\left(C_1\cos(\beta t)+C_2\sin(\beta t)\right).
$$
Can we work in reverse to find the differential equation? This solution is derived from a characteristic equation with complex eigenvalues
$$
\lambda=\alpha \pm i\beta.
$$
That means we can recover the quadratic equation by doing the steps of completing the square but in reverse. Let's do an example.

\begin{example}[3]
{
Find a differential equation that has the solution 
$$
x = e^{-2t}(3\cos(4t)-\sin(4t)).
$$
}
First, we identify the eigenvalues. Since $\alpha = -2$ and $\beta=4$, the eigenvalues are 
$$
\lambda = -2\pm 4i.
$$
Then we perform the steps of completing the square in reverse 
\begin{align*}
    \lambda &= -2\pm 4i \\
    \lambda +2 &= \pm 4i && \com{Add 2}\\
    (\lambda +2)^2 &= (\pm 4)^2(i)^2 && \com{Square both sides}\\
    \lambda^2 +4\lambda + 4 &= -16   && \com{Simplify}\\
    \lambda^2 +4\lambda + 20 &= 0 && \com{Add 16}
\end{align*}
Finally, we take the coefficients of the characteristics $(a=1,b=4,c=20)$ and write the differential equation
$$
\boxed{
    x'' +4x'+20x=0.
}
$$
\end{example}
There are other ways to find the characteristic equation. For example, you may write
\[
(\lambda-(-2+4i))(\lambda-(-2-4i))=0.
\]
and expand it. It will eventually give you same equation $\lambda^2+4\lambda+20=0$. You can also use the quadratic formula in reverse. Alternatively, choose $a=1$ and compare
\[
\lambda=-\frac{b}{2}\pm\frac{\sqrt{4c-b^2}}{2}i
\]
with $\lambda=-2\pm4i$. Then
\[
-\frac{b}{2}=-2 \implies b=4,
\qquad
\frac{\sqrt{4c-16}}{2}=4 \implies c=20.
\]
You have three options to choose from. Choose whichever you feel comfortable with.

% ---------------------------------

\subsection{Shape of the Solution}

Let's briefly discuss how the solution appears when plotted on the $tx$-plane.
A complex-root solution is the product of an exponential and a wave:
\[
x(t)=
\underbrace{e^{\alpha t}}_{\text{exponential}}
\underbrace{(C_1\cos\beta t+C_2\sin\beta t)}_{\text{wave}}.
\]
The parameter $\beta>0$ represents the frequency which controls how fast the wave oscillates.
The parameter $\alpha$ determines its amplitude: growing when
$\alpha>0$, decaying when $\alpha<0$, and remaining constant when $\alpha=0$. The sign of $\alpha$ determines the three main differences in the behavior of the solution. The three cases are shown in the figure below.

\begin{figure}[H]
\centering

% Growing wave
\begin{subfigure}[t]{0.32\textwidth}
\centering
\resizebox{\linewidth}{!}{%
\begin{tikzpicture}
    \draw[->,thick] (-0.3,0) -- (12.2,0);
    \draw[->,thick] (0,-3.5) -- (0,3.5);

    \draw[thick,domain=0:11.8,samples=350]
        plot (\x,{exp(0.1*\x)*cos(60*\x)});
    \draw[dashed,domain=0:11.8,samples=150]
        plot (\x,{exp(0.1*\x)});
    \draw[dashed,domain=0:11.8,samples=150]
        plot (\x,{-exp(0.1*\x)});
\end{tikzpicture}%
}
\caption{Growing wave\\$\alpha>0$.}
\end{subfigure}
\hfill
% Decaying wave
\begin{subfigure}[t]{0.32\textwidth}
\centering
\resizebox{\linewidth}{!}{%
\begin{tikzpicture}
    \draw[->,thick] (-0.3,0) -- (12.2,0);
    \draw[->,thick] (0,-3.5) -- (0,3.5);

    \draw[thick,domain=0:11.8,samples=350]
        plot (\x,{2*exp(-0.1*\x)*cos(60*\x)});
    \draw[dashed,domain=0:11.8,samples=150]
        plot (\x,{2*exp(-0.1*\x)});
    \draw[dashed,domain=0:11.8,samples=150]
        plot (\x,{-2*exp(-0.1*\x)});
\end{tikzpicture}%
}
\caption{Decaying wave\\$\alpha<0$.}
\end{subfigure}
\hfill
% Pure oscillation
\begin{subfigure}[t]{0.32\textwidth}
\centering
\resizebox{\linewidth}{!}{%
\begin{tikzpicture}
    \draw[->,thick] (-0.3,0) -- (12.2,0);
    \draw[->,thick] (0,-3.5) -- (0,3.5);

    \draw[thick,domain=0:11.8,samples=350]
        plot (\x,{2*cos(60*\x)});
    \draw[dashed] (0,2) -- (11.8,2);
    \draw[dashed] (0,-2) -- (11.8,-2);
\end{tikzpicture}%
}
\caption{Pure oscillation\\$\alpha=0$.}
\end{subfigure}

\caption{The real part of the eigenvalue determines whether
the amplitude grows, decays, or stays constant.
Dashed curves indicate the envelopes.}
\end{figure}

%%%%%%%%%%%%%%%%%%%%%%%%%%%%%%%%%%%%%%%%%%%%%%%%%%%%%%%%%%%%%%%%%%%%%%%%
\section{Summary of the Characteristic Equation Method}


\subsection{Table}
Over Days 12–14, we have learned to solve
\[
ax''+bx'+cx=0
\]
by finding the roots of the characteristic equation
\[
a\lambda^2+b\lambda+c=0.
\]
The sign of the discriminant $D=b^2-4ac$ determines which solution form to use. The three cases are summarized in the table below. It is useful to memorize the general solutions.

\begin{table}[H]
\centering
\small
\renewcommand{\arraystretch}{1.0}
\setlength{\tabcolsep}{4pt}
\begin{tabular}{|
    >{\centering\arraybackslash}m{2.0cm}|
    >{\centering\arraybackslash}m{2.0cm}|
    >{\centering\arraybackslash}m{2.3cm}|
    >{\centering\arraybackslash}m{1.7cm}|
    >{\centering\arraybackslash}m{4.5cm}|
}
\hline
\textbf{Roots}
& \shortstack{\textbf{General}\\\textbf{eigenvalues}}
& \textbf{Example DE}
& \shortstack{\textbf{Example}\\\textbf{eigenvalues}}
& \textbf{General solution} $x(t)$
\\
\hline
\shortstack{Real, distinct\\$D>0$}
&
$\lambda_1\ne\lambda_2$
&
$x''-x'-2x=0$
&
\shortstack{$\lambda_1=2$\\$\lambda_2=-1$}
&
$C_1e^{\lambda_1t}+C_2e^{\lambda_2t}$
\\
\hline
\shortstack{Real, repeated\\$D=0$}
&
$\lambda$
&
$x''-4x'+4x=0$
&
$\lambda=2$
&
$(C_1+C_2t)e^{\lambda t}$
\\
\hline
\shortstack{Complex\\$D<0$}
&
$\lambda=\alpha\pm i\beta$
&
$x''-2x'+5x=0$
&
$\lambda=1\pm2i$
&
$C_1e^{\alpha t}\cos(\beta t)+C_2e^{\alpha t}\sin(\beta t)$
\\
\hline
\end{tabular}
\caption{Summary of the cases for the second-order linear constant coefficient DEs}
\end{table}

\subsection{Tips on Solving the Quadratic}

The characteristic equation is a quadratic, so we can solve it using familiar algebraic techniques:
\begin{itemize}
\item \textbf{Factoring:} Try factoring when the factors are
easy to recognize. Try to rewrite the quadratic in the form
$(\lambda-r_1)(\lambda-r_2)$. If the factors are not apparent,
use another method. Here is an example,
    \[
    \lambda^2-\lambda-2
    =(\lambda-2)(\lambda+1)=0
    \quad\Rightarrow\quad \lambda=2,-1.
    \]
    If you want to be safe, you can verify your factorization by FOILing.
    
    \item \textbf{Completing the square:} If you know for sure that the roots are complex, then try to complete the square.
    \[
    \lambda^2+2\lambda+10=0
    \quad\Rightarrow\quad
    (\lambda+1)^2=-9
    \quad\Rightarrow\quad \lambda=-1\pm3i.
    \]

    \item \textbf{Quadratic formula:} This method always works, even when
    the factors are not obvious. You should memorize this. 
    \[
    \lambda=\frac{-b\pm\sqrt{b^2-4ac}}{2a}.
    \]
\end{itemize}
Last tip, try to simplify the equation when you spot a common factor, or have fractions in the characteristic equation. For example,
\begin{align*}
3x''+9x'+27x=0
&\quad\implies\quad x''+3x'+9x=0,\\
\tfrac12 x''+8x'+4x=0
&\quad\implies\quad x''+16x'+8x=0.
\end{align*}

%%%%%%%%%%%%%%%%%%%%%%%%%%%%%%%%%%%%%%%%%%%%%%%%%%%%%%%%%%%%%%%%%%%%%%%%
\section{Connection to Springs}

\subsection{Review Spring Model}

Let $x(t)$ be the displacement from equilibrium of a ball attached
to a spring moving over a table. Combining Newton's second law,
Hooke's law, and a damping force gives \textbf{spring-mass equation}
\begin{align*}
\sum F&=mx'',\\
-kx-\gamma x'&=mx'',\\
mx''+\gamma x'+kx&=0.
\end{align*}
Here, $m>0$ is the mass, $\gamma\ge0$ is the damping coefficient,
and $k>0$ is the spring constant. This is exactly the type of
second-order linear constant-coefficient DE we have been solving.


%-------------------------------------------------------------
\subsection{Key Differences}

There is one key difference between
\[
ax''+bx'+cx=0 \quad \text{versus}\quad mx''+\gamma x'+kx=0.
\]
In the general equation, the coefficients are real and $a\ne0$,
but their signs are otherwise unrestricted. In the spring model,
the physical parameters must satisfy
\[
m>0,\qquad \gamma\ge0,\qquad k>0.
\]
The mass and spring constant are positive, and damping opposes
the motion. The case $\gamma=0$ represents a spring with no damping.

These restrictions affect the possible solutions. For example,
a general constant-coefficient DE may have a growing solution such as
\[
y(t)=e^{2t}\cos3t.
\]
If $y(t)$ describes the displacement of the mass, its amplitude
increases with each swing. Such behavior cannot occur in our
unforced spring-mass model: the spring restores the mass toward
equilibrium, while damping removes energy rather than supplying it.

It turns out that a solution such as $y(t)=e^{2t}\cos3t$ is \textit{impossible} to obtain as a solution to 
$$
mx''+\gamma x'+kx=0.
$$
In a later lecture, we will explain why this is so. We will connect all three root types to
different kinds of spring motion.

\end{document}